\frontchapter{Abstract}
As AI model training scales across growing numbers of GPUs, communication becomes a bottleneck. Most communication is host-orchestrated: launches and synchronization can leave devices waiting. GPU invocation moves control into kernels, while only some transports additionally let the GPU post network work directly to the NIC. Neither property is inherently faster; pattern, size, topology, synchronization, and progress determine the result. This thesis characterizes complete communication paths and tests which observations remain valid in a Conjugate Gradient (CG) solver on the Leonardo Booster partition.

We develop \emph{gpu-comm-benchmark}, a CUDA/SYCL suite comparing MPI, NCCL/oneCCL, NVSHMEM, and OSHMPI across point-to-point, halo, collective, and CG-inspired patterns. We also extend oneCCL with grouped NCCL point-to-point, an OSHMPI backend for the benchmarked world-communicator operation set, and a validated NVSHMEM runtime and staging foundation, and extend aCG with OSHMPI and SYCL implementations.

Experiments on 1--32 GPUs expose a decision boundary rather than one winner. MPI minimizes startup cost; persistent, kernel-invoked NVSHMEM leads large steady multi-node halos; NCCL paths lead large allreduce and multi-node all-to-all. MPI/UCC avoids a 17--34$\times$ small-message all-to-all slowdown at 16 and 32 ranks, while its larger multi-node allreduce configurations benefit above a size-dependent crossover around 64\,KiB. In the reconciled synthetic \texttt{cg\_step} archive, OSHMPI has the lowest point estimate in 25 of 30 multi-rank cells, including every multi-node cell, but its archived scalar result remains in host memory; NCCL leads the remaining five and the single-GPU controls.

Comparing \texttt{cg\_step} with aCG gives 7.8--11.1\% mean absolute discrepancy for the matched MPI and NCCL scalar-reduction intervals. The regular synthetic halo misses the application's irregular peers and overlap, and the archived OSHMPI paths also differ in scalar primitive and result placement. NCCL has the lowest representative host-initiated solver times. Device-side NVSHMEM uses CPU-proxy-assisted inter-node networking and improves relative to blocking OSHMPI at larger rank counts but remains behind NCCL. The evidence therefore supports configured performance regimes and selective primitive-level agreement, not an initiation-only speedup; complete application selection requires matching communication structure and completion dependencies.

\printkeywords
